\documentclass[12pt, a4paper]{article}

\usepackage[utf8]{inputenc}
\usepackage[T1]{fontenc}
\usepackage{amsmath, amssymb}
\usepackage{graphicx}
\usepackage{booktabs}
\usepackage{hyperref}
\usepackage[margin=2.5cm]{geometry}
\usepackage{natbib}
\usepackage{xcolor}
\usepackage{float}

\bibliographystyle{plainnat}

\newcommand{\tbd}[1]{\textcolor{red}{\textbf{[TBD: #1]}}}

\title{Explainable EEG ``Brain-Age Clock'' in Childhood--Adolescence\\Using Healthy Brain Network Data\\[0.5em]\large Pilot Study}
\author{Zoe Research Team}
\date{\today}

\begin{document}

\maketitle

\begin{abstract}
Electroencephalography (EEG) provides a non-invasive window into brain maturation, yet most EEG-based age-prediction studies treat the model as a black box or conflate periodic (oscillatory) and aperiodic (1/$f$ background) spectral components.
We present a reproducible, open-source pipeline for building an explainable ``brain-age clock'' from resting-state EEG, with explicit separation of periodic band powers, peak alpha frequency (PAF), individualized alpha power, and aperiodic spectral parameters via spectral parameterization (specparam).
In this pilot study, we demonstrate the full pipeline on a 20-participant subset of a public OpenNeuro dataset (ds004148, ages 18--22), extracting 34 spectral features per recording and evaluating age prediction with leakage-safe cross-validation.
Pilot results show that the narrow age range limits predictive performance (best MAE = 1.01 years, baseline), but the pipeline successfully extracts physiologically meaningful features (mean PAF = 10.2~Hz, aperiodic exponent = 0.91) and identifies parietal theta power as the most informative feature.
We describe plans for scaling to the full Healthy Brain Network dataset (ages 5--21, $N > 2{,}000$) where developmental EEG trajectories are expected to provide substantially greater predictive signal.
\end{abstract}

\vspace{1em}
\noindent\textbf{Guardrails:} This work is strictly for research and education. The pipeline outputs are not diagnostic tools, clinical decision-support systems, or validated biomarkers. Model predictions have not been clinically validated.

\section{Introduction}

The human brain undergoes profound structural and functional changes from childhood through adolescence. These maturational processes---including synaptic pruning, myelination, and the refinement of cortical circuits---produce measurable changes in the brain's electrical activity that can be captured non-invasively through electroencephalography (EEG).

EEG signals can be decomposed into spectral components reflecting distinct neural processes. Oscillatory (``periodic'') components, such as the alpha rhythm ($\sim$8--13~Hz), arise from synchronized neural populations and reflect thalamocortical circuit function. The frequency at which alpha power peaks---the peak alpha frequency (PAF)---increases systematically during childhood and adolescence, from roughly 6--8~Hz in early childhood to 9--11~Hz in late adolescence \citep{miskovic2015development, cellier2021alpha}.

However, the EEG power spectrum also contains a broadband, aperiodic ``1/$f$-like'' background whose slope (exponent) and level (offset) change with development: younger children show steeper spectral slopes that flatten during adolescence \citep{donoghue2020specparam, ostlund2022specparam}. If periodic and aperiodic components are not separated, apparent changes in ``alpha power'' may actually reflect shifts in the aperiodic background.

The specparam algorithm \citep{donoghue2020specparam} enables explicit decomposition of the power spectrum into periodic peaks and aperiodic background. Applying this decomposition to developmental EEG offers more physiologically precise features and more interpretable age-prediction models.

In this pilot study, we present a complete, reproducible pipeline for building an explainable EEG brain-age clock with explicit periodic/aperiodic separation. We demonstrate it on a small dataset and describe plans for scaling to the Healthy Brain Network (HBN).

\section{Related Work}

\subsection{EEG-BIDS and Reproducible Research}
The Brain Imaging Data Structure for EEG (EEG-BIDS; \citealt{pernet2019eegbids}) standardizes how EEG data are organized, enabling reproducible large-scale analyses. The OpenNeuro platform \citep{markiewicz2021openneuro} hosts BIDS-formatted datasets. Our pipeline uses MNE-BIDS \citep{appelhoff2019mnebids} to read data directly from the BIDS layout.

\subsection{The Healthy Brain Network}
The Healthy Brain Network (HBN; \citealt{alexander2017hbn}) includes EEG recordings from children and adolescents aged approximately 5--21, providing one of the largest publicly available pediatric EEG resources.

\subsection{Developmental EEG Trajectories}
PAF increases from $\sim$6--8~Hz in early childhood to 9--11~Hz by late adolescence \citep{miskovic2015development, cellier2021alpha}. Delta and theta power decrease with age \citep{clarke2001age, whitford2007maturation}. The aperiodic exponent flattens during development \citep{ostlund2022specparam}.

\subsection{Spectral Parameterization}
The specparam algorithm \citep{donoghue2020specparam} models the power spectrum as the sum of an aperiodic component and Gaussian peaks, enabling separate analysis of oscillatory and background contributions.

\subsection{EEG-Based Brain-Age Models}
Brain-age prediction from EEG has been demonstrated in adults \citep{engemann2022brainage, sun2019brainage}, but most prior work has focused on adult populations, used the model as a black box, or did not explicitly separate periodic and aperiodic components.

\section{Data}

\subsection{Pilot Dataset}
We used a publicly available resting-state EEG dataset (OpenNeuro ds004148), comprising recordings from healthy adults with eyes-open (EO) and eyes-closed (EC) conditions across multiple sessions. We downloaded data from 20 participants (seeded random sample), yielding 120 EEG recordings. Participants ranged in age from 18 to 22 years (mean = 19.8).

\textbf{Note:} The intended target is the HBN dataset (ds004186, ages 5--21). The pilot dataset was used because ds004186 requires authentication. \tbd{run on full HBN dataset}

\subsection{Ethics}
All data are publicly available, de-identified, and obtained under OpenNeuro data use agreements. This project is strictly for research and education purposes.

\section{Methods}

\subsection{Quality Control and Preprocessing}
Minimal preprocessing was applied: bandpass filter (1--40~Hz), notch filter (60~Hz), common average re-reference, 2-second epoching, and amplitude-based artifact rejection (peak-to-peak $>$~200~$\mu$V or flatline $<$~1~$\mu$V).

\subsection{Feature Engineering}
From each cleaned recording, 34 spectral features were extracted in four families:

\begin{itemize}
    \item \textbf{Band powers (24):} Global and regional (frontal, central, parietal, occipital, temporal) log$_{10}$ absolute power in delta (1--4~Hz), theta (4--8~Hz), alpha (8--12~Hz), and beta (13--30~Hz).
    \item \textbf{PAF and individualized alpha (4):} Peak alpha frequency from posterior channels (6--13~Hz search range), individualized alpha power (PAF~$\pm$~2~Hz), and specparam alpha peak parameters.
    \item \textbf{Aperiodic (2):} Exponent and offset from specparam fit (2--40~Hz range).
    \item \textbf{Ratios (2):} Theta/alpha and theta/beta power ratios.
\end{itemize}

\subsection{Modeling}
Three models were evaluated: baseline (mean), RidgeCV, and HistGradientBoosting. GroupKFold CV ($K=5$) grouped by participant ID prevented data leakage.

\subsection{Explainability}
Permutation importance was computed on held-out folds. Ablation experiments removed PAF, aperiodic, or both feature groups.

\section{Results (Pilot)}

\subsection{Cohort Flow}

\begin{table}[H]
\centering
\caption{Pilot cohort flow}
\begin{tabular}{lr}
\toprule
Stage & Recordings \\
\midrule
Downloaded & 120 \\
Passed QC & 120 (100\%) \\
Feature extraction & 118 (98.3\%) \\
Included in modeling & 118 \\
\bottomrule
\end{tabular}
\label{tab:cohort}
\end{table}

\subsection{Model Performance}

\begin{table}[H]
\centering
\caption{Pilot model performance (5-fold GroupKFold CV)}
\begin{tabular}{lcc}
\toprule
Model & MAE (years) & $R^2$ \\
\midrule
Baseline (mean) & 1.01 $\pm$ 0.43 & $-$0.115 $\pm$ 0.060 \\
RidgeCV & 1.25 $\pm$ 0.42 & $-$0.997 $\pm$ 0.610 \\
HGB & 1.12 $\pm$ 0.48 & $-$0.589 $\pm$ 0.310 \\
\bottomrule
\end{tabular}
\label{tab:performance}
\end{table}

All $R^2$ values are negative, indicating that no model outperformed predicting the mean---expected given the 4-year age range.

\subsection{Feature Importance}

\begin{figure}[H]
\centering
\includegraphics[width=0.9\textwidth]{figures/fig5_importance.png}
\caption{Top 20 features by permutation importance (HGB model). Parietal theta power dominates.}
\label{fig:importance}
\end{figure}

\subsection{Feature Distributions}

\begin{figure}[H]
\centering
\begin{minipage}{0.48\textwidth}
\includegraphics[width=\textwidth]{figures/fig3_age_vs_paf.png}
\end{minipage}
\hfill
\begin{minipage}{0.48\textwidth}
\includegraphics[width=\textwidth]{figures/fig3_aperiodic_hist.png}
\end{minipage}
\caption{Left: Age vs.\ peak alpha frequency (pilot). Right: Aperiodic exponent distribution. Mean PAF = 10.2~Hz; mean exponent = 0.91.}
\label{fig:features}
\end{figure}

\subsection{Ablation Results}

\begin{table}[H]
\centering
\caption{Ablation study (HGB model, pilot data)}
\begin{tabular}{lcc}
\toprule
Experiment & MAE (years) & Features \\
\midrule
Full model & 1.12 & 34 \\
Remove PAF/IAF & 1.13 & 30 \\
Remove aperiodic & 1.15 & 32 \\
Remove both & 1.10 & 28 \\
\bottomrule
\end{tabular}
\label{tab:ablation}
\end{table}

\subsection{Predicted vs.\ True Age}

\begin{figure}[H]
\centering
\begin{minipage}{0.48\textwidth}
\includegraphics[width=\textwidth]{figures/fig4_pred_vs_true.png}
\end{minipage}
\hfill
\begin{minipage}{0.48\textwidth}
\includegraphics[width=\textwidth]{figures/fig4_residuals.png}
\end{minipage}
\caption{Left: Predicted vs.\ true age (HGB, CV predictions). Right: Residuals vs.\ age.}
\label{fig:calibration}
\end{figure}

\section{Discussion}

The pilot results confirm that the pipeline operates correctly end-to-end. However, the narrow age range (18--22) fundamentally limits developmental signal. Parietal theta power emerged as the most important feature, likely reflecting individual differences in attentional state rather than maturation. Ablation differences were minimal ($\leq$0.03 years).

We expect that applying this pipeline to the full HBN dataset (ages 5--21) will reveal increasing PAF, flattening aperiodic slopes, and shifting regional power distributions, providing substantially more predictive signal.

\section{Limitations}

\begin{enumerate}
    \item \textbf{Narrow age range:} The pilot spans only 18--22 years, far narrower than the target 5--21.
    \item \textbf{Small sample:} 20 participants yield high CV variance.
    \item \textbf{Proxy dataset:} ds004148 differs from the target HBN dataset.
    \item \textbf{No ICA:} Only amplitude-based artifact rejection was applied.
    \item \textbf{Not diagnostic:} This pipeline is a research tool; it has not been clinically validated.
\end{enumerate}

\section{Next Steps}

\tbd{Obtain OpenNeuro authentication for ds004186 (HBN). Target: $N > 500$ participants, ages 5--21.}

Key upgrades for the full study:
\begin{itemize}
    \item Apply pipeline to full HBN dataset across the developmental age range
    \item Compare EO-only, EC-only, and combined models
    \item Test preprocessing sensitivity (filter bounds, reference choice)
    \item Add external validation on a second dataset
\end{itemize}

\section{Reproducibility}

\begin{itemize}
    \item Dataset: OpenNeuro ds004148 (pilot), ds004186 (target)
    \item Random seed: 1337
    \item Software: Python 3.13, MNE 1.12, scikit-learn 1.7, specparam
    \item All stages are scripted and config-driven; re-run stages 0--6 sequentially
\end{itemize}

\bibliography{references}

\end{document}
